\begin{abstract}
Fraud rings share devices, cards and network addresses and move money among
themselves, which graph queries can expose. Two obstacles appear on live
transaction streams. Legitimate shared infrastructure (office gateways, carrier
NAT, kiosks) links many unrelated accounts, and the usual remedy, a hard degree
cap, hides any ring that shares a busy attribute. Whole-graph queries re-run on
a timer also cost more as the graph grows. We replace the cap with a continuous
hub-aware weight with time decay, detect incrementally on the $k$-hop region
touched by each micro-batch, and pass candidate rings through a calibrated
scorer that explains every alert. On a deterministic simulator with ring-level
ground truth and five evasion strategies, over \EScenarios{} scenarios and
\ESeeds{} test seeds, the method reaches ring F1 \EPropF{} against \ECapTenF{}
for the original degree-capped Cypher rules and \ELouvainF{} for Louvain
communities (paired Wilcoxon $p<0.001$). Hub weighting extends the reachable
recall on shared-attribute rings from \FrCapSMaxR{} to \FrHubSMaxR{}, though
not at higher precision everywhere, and the scorer contributes most of the
precision. Incremental detection flags the same rings as full recomputation
and at \ScaleMaxUsers{} users takes \ScaleIncMs{}\,ms per batch against
\ScaleFullMs{}\,ms. Rings that rotate identities or act slowly evade every
method we tested. All results are synthetic; we state what the public AML
datasets we could use can and cannot test.
\end{abstract}
